% Bloque B: Pedro, 2:45. S6--S8. El problema: tormenta de reconexiones,
% falla metaestable y qué evidencia hay (y no hay) en Binance.
\section{El problema}
\begin{frame}{El problema: una tormenta de reconexiones}
  \centering% S6. El bucle de la tormenta de reconexiones.
\begin{tikzpicture}[x=1cm,y=1cm]
  \path[use as bounding box] (0,0.1) rectangle (13.5,4.65);
  \node[malo,font=\scriptsize,text width=2.8cm] (t) at (1.7,3.95) {\textbf{Disparador}\\pico de volatilidad\\con 10$\times$ más mensajes};
  \node[box,font=\scriptsize,text width=3.7cm] (n1) at (6.6,3.95) {\textbf{1.} Clientes lentos no alcanzan\\$\to$ se les \textbf{corta} la conexión};
  \node[box,font=\scriptsize,text width=3.4cm] (n2) at (10.9,2.3) {\textbf{2.} Miles piden el\\\textbf{snapshot completo}\\por REST \textbf{a la vez}};
  \node[box,font=\scriptsize,text width=3.7cm] (n3) at (6.6,0.65) {\textbf{3.} El servidor central se satura\\(stream 10$\times$ + snapshots)};
  \draw[flow,colWake] (t) -- (n1);
  \draw[flow] (n1.east) to[out=0,in=90] (n2.north);
  \draw[flow] (n2.south) to[out=270,in=0] (n3.east);
  \draw[flow,colWake,line width=1.2pt] (n3.west) to[out=180,in=180,looseness=1.6]
    node[pos=0.5,left,font=\scriptsize,text=colWake,align=right] {el stream se retrasa\\y más clientes\\se atrasan} (n1.west);
  \node[font=\Large\bfseries,text=colWake] at (7.6,2.3) {bucle};
  \node[malo,font=\scriptsize,text width=2.8cm] at (1.7,1.0) {Los retries repetidos\\llevan a \textbf{429} $\to$ \textbf{418} (ban)};
\end{tikzpicture}

  \takeaway{La recuperación compite con el stream por la misma capacidad.}
  \note{Pedro. 4:20--5:20 (60 s). Ahora imaginen un momento de mucha
  volatilidad: el bitcoin cae fuerte y el número de mensajes se multiplica,
  en nuestro experimento por diez. Los clientes con internet lento no
  alcanzan a recibir todo y se les corta la conexión. Según la receta de
  Binance, cada uno pide el libro completo, y son miles a la vez. El
  servidor central tiene que mandar el stream, que ya va a diez veces, y
  además miles de snapshots. Se satura y retrasa el stream, y eso atrasa a
  más clientes, que también se cortan y piden snapshot. Es un bucle. Y los
  límites lo empeoran: quien reintenta mucho termina baneado.}
\end{frame}

\begin{frame}{¿Por qué es grave? La falla metaestable}
  \centering% S7. Sobrecarga normal vs falla metaestable. Esquema cualitativo, no datos.
\begin{tikzpicture}[x=1cm,y=1cm]
  \path[use as bounding box] (0,-0.35) rectangle (13.5,4.5);
  \foreach \dx/\titulo/\col in {0.9/Sobrecarga normal/colReq,7.9/Falla metaestable/colWake} {
    \begin{scope}[xshift=\dx cm]
      \node[anchor=west,font=\small\bfseries,text=\col] at (0,4.25) {\titulo};
      \fill[colWake!9] (1.5,0) rectangle (2.9,3.6);
      \node[font=\scriptsize,text=colWake,align=center] at (2.2,3.3) {pico\\(disparador)};
      \draw[flow,thin] (0,0) -- (5.6,0);
      \draw[flow,thin] (0,0) -- (0,3.75);
      \node[anchor=north east,font=\tiny] at (5.6,-0.05) {tiempo};
      \node[rotate=90,font=\tiny,anchor=south] at (-0.1,1.35) {clientes con libro válido};
      \draw[colNeutro!50,densely dotted] (0,2.8) -- (5.4,2.8);
      \node[font=\tiny,anchor=east] at (0,2.8) {100\%};
    \end{scope}
  }
  % Sobrecarga normal: cae durante el pico y vuelve sola
  \begin{scope}[xshift=0.9cm]
    \draw[colReq,line width=1.6pt] (0,2.8) -- (1.5,2.8) -- (1.75,1.0) -- (2.9,1.0) -- (3.3,2.8) -- (5.4,2.8);
    \node[font=\scriptsize,text=colReq,align=center] at (4.3,2.2) {se recupera\\sola};
  \end{scope}
  % Metaestable: el pico termina y el sistema sigue abajo
  \begin{scope}[xshift=7.9cm]
    \draw[colWake,line width=1.6pt] (0,2.8) -- (1.5,2.8) -- (1.75,0.6) -- (2.9,0.6) -- (3.4,0.8) -- (4.2,0.55) -- (5.4,0.7);
    \node[font=\scriptsize,text=colWake,align=center] at (4.2,1.7) {sigue caída:\\los reintentos\\la sostienen};
  \end{scope}
\end{tikzpicture}

  \takeaway{Metaestable: el pico ya pasó, pero el sistema sigue caído.}
  \source{Esquema cualitativo, no son mediciones. Isaacs et al., \emph{Analyzing Metastable Failures}, HotOS 2025.}
  \note{Pedro. 5:20--6:20 (60 s). Este bucle tiene nombre en la literatura:
  falla metaestable. A la izquierda, una sobrecarga normal: durante el pico
  los clientes pierden su libro, pero cuando el pico termina el sistema se
  recupera solo. A la derecha, una falla metaestable: el pico termina, la
  carga vuelve a lo normal, y el sistema sigue caído, porque el efecto
  sostenedor, en este caso los reintentos de snapshots, lo mantiene
  saturado. Suele hacer falta intervención manual: cortar tráfico o
  reiniciar. Una analogía: un accidente provoca un atasco; retiran el auto,
  pero el atasco sigue, porque los autos que siguen llegando lo sostienen.
  La referencia principal es Isaacs y otros, 2025, que propone analizar
  estas fallas con modelos y simulación. Eso es lo que hicimos.}
\end{frame}

\begin{frame}{¿Le pasa esto a Binance de verdad?}
  \centering% S8. Qué evidencia hay y qué no hay de que Binance sufra este bucle.
\begin{tikzpicture}[x=1cm,y=1cm]
  \path[use as bounding box] (0,0.1) rectangle (13.5,4.6);
  \tikzset{item/.style={box,font=\scriptsize,text width=5.7cm,align=left,anchor=west}}
  \node[anchor=west,font=\small\bfseries] at (0.1,4.35) {Lo que sí hay: indicios};
  \node[item] at (0.1,3.45) {Usuarios del \textbf{foro oficial} de Binance reportan conexiones cerradas con el código \textbf{1006} durante trading intenso.};
  \node[item] at (0.1,2.3) {También reportan \textbf{muchos streams de profundidad} cerrándose al mismo tiempo.};
  \node[item] at (0.1,1.15) {\textbf{Tardis.dev}, proveedor de datos, reconstruye el libro con el \textbf{snapshot REST} cuando pierde un mensaje.};
  \draw[colNeutro!40] (6.75,0.15) -- (6.75,4.55);
  \node[anchor=west,font=\small\bfseries,text=colWake] at (7.1,4.35) {Lo que no hay};
  \node[item,malo,align=left] at (7.1,3.45) {Ningún \textbf{post-mortem público} de Binance atribuye una caída a este bucle.};
  \node[item,malo,align=left] at (7.1,2.3) {No conocemos sus \textbf{mitigaciones internas} ni su capacidad real.};
  \node[item,malo,align=left] at (7.1,1.15) {No sabemos \textbf{cuánto le cuesta generar un snapshot}: lo variamos en la simulación.};
\end{tikzpicture}

  \takeaway{Afirmamos que el contrato \emph{tiene la estructura} del problema, no que Binance ya falló.}
  \note{Pedro. 6:20--7:05 (45 s). Tenemos que ser honestos. Hay indicios:
  en el foro oficial de desarrolladores de Binance hay usuarios que
  reportan conexiones cerradas con el código 1006 en momentos de mucho
  trading, y muchos streams de profundidad cerrándose a la vez. Y
  Tardis.dev, un proveedor de datos históricos, documenta que reconstruye
  el libro con el snapshot REST cuando pierde un mensaje. Pero no hay
  ningún post-mortem público de Binance que atribuya una caída a este
  bucle, y no conocemos sus mitigaciones internas ni cuánto le cuesta
  generar un snapshot. Por eso solo afirmamos que el contrato público
  cumple las condiciones de una falla metaestable, y lo ponemos a prueba
  con una simulación. Nicolás presenta nuestra propuesta.}
\end{frame}
